\newpage
\thispagestyle{empty}
\begin{center}
  \textbf{\MakeUppercase{Abstract}}
\end{center}

\noindent
This dissertation investigates how optimization problems can be understood
within a unified compositional and diagrammatic framework grounded in category
theory. The central observation is that key mathematical structures in
optimization can be interpreted as morphisms in categories whose
composition corresponds to the elimination or minimization of internal
variables. From this perspective, optimization emerges not as a collection of
isolated techniques but as an algebraic operation intrinsic to relational
structure.

The development follows an incremental progression rooted in Lawvere quantales
and weighted relations, adding one new generator and a finite set of axioms at
each step. We recall the theory of Quantale relations 
and the Polyhedral Algebra framework as preliminaries steps
to develop the novel Graphical Linear Programming (GLP),
where morphisms -- diagrams -- are parametric linear programs. Those composes
with other optimization blocks, so that the elimination of internal variables,
or, in this case, the optimization itself, corresponds exactly to categorical composition.
Soundness and completeness results for the previous algebras are recalled from prior works,
and in the last section we also present the completeness theorem for our new language.

The diagrammatic calculus plays the same role for
linear programming that Graphical Linear Algebra plays for matrices: it provides
a expressible language for normal forms and rewriting rules in
which arguments can be carried out symbolically, prior to and
independently of numerical computation, making the main properties of the
system evident in the syntax itself. The tools and results developed and proved
here are this diagrammatic calculus for Linear Programming.

\vspace{1em}
\noindent
Keywords: category theory; string diagrams; linear programming; convex
optimization; symmetric monoidal categories; weighted relations; quantales;
graphical algebra; normal forms; completeness.
\clearpage
